\documentclass[11pt]{article}

\usepackage[margin=1in]{geometry}
\usepackage{amsmath,amssymb,amsthm}
\usepackage[round,authoryear]{natbib}
\usepackage{microtype}
\usepackage{booktabs,graphicx}
\usepackage[hidelinks]{hyperref}

% amsthm provides the proof environment but custom theorem-like
% environments must be declared explicitly before \begin{document}.
\newtheorem{proposition}{Proposition}[section]

% cleveref must load after hyperref and theorem definitions.
\usepackage[noabbrev]{cleveref}

\begin{document}

\title{Horizon-Matched Cross-Validation for Forecast-Optimal Penalized Trend Smoothness}

\author{Heriberto Espino Montelongo\\
Universidad de las Am\'ericas Puebla, Puebla, Mexico}
\date{}

\maketitle

\begin{abstract}
Penalized least-squares trend estimation requires a smoothing
parameter, but a good historical fit need not extrapolate well.
We investigate chronological cross-validation for selecting
smoothness according to its future predictive performance.
A normalized trace-based coordinate $S\in[0,1]$ represents
the complete range of penalties, including the unsmoothed fit
and the limiting polynomial trend. For each candidate $S$,
the smoothing matrix is applied to a historical window,
the fitted trend is continued over the intended $h$-step
forecast horizon, and the subsequent block is scored by
mean squared error. The selected smoothness minimizes pooled
historical forecast loss and is refitted at the operational
origin without using future test observations. The PLS
estimator, smoothness index, and forecast-oriented tuning
have established precedents; the object studied here is
their particular horizon-specific combination for
finite-difference trend extrapolation. A frozen simulation
with 3,000 scenarios and 72,000 outer-origin--horizon
evaluations compares the procedure with one-step forecast
tuning and conventional criteria. At horizons $3$, $6$,
and $12$, geometric forecast-RMSFE ratios relative to
one-step tuning are $0.947$, $0.861$, and $0.762$.
The criterion can be multimodal, requiring comparison
of interior candidates and the limiting endpoints.
\end{abstract}

\noindent\textbf{Keywords:} time-series cross-validation; penalized least squares; forecast mean squared error; normalized smoothness; trend extrapolation; Monte Carlo simulation.

\medskip

\section{Introduction}
\label{sec:introduction}

A trend smoother that reconstructs historical observations well
need not produce an accurate forecast. In finite-difference
penalized least squares (PLS), the smoothing penalty trades
fidelity to the past against regularity of the fitted path.
Once the path is extrapolated, however, its terminal
level and differences determine the forecast. Selecting
the penalty to recover a historical trend and selecting
it for future extrapolation are different problems.

This paper investigates a specific predictive criterion.
For a historical window $x_t$, let $H(S)$ be the PLS
smoothing matrix indexed by normalized smoothness
$S\in[0,1]$, and let $G_{d,h}$ continue the fitted
trend $h$ observations into the future. At the final
forecast origin $T$, choose
\begin{equation}
\label{eq:intro_pooled_loss}
F^{\mathrm{pool}}_{T,d,L,h}(S)
=\frac{1}{Mh}\sum_{m=1}^M
\left\|y_{t_m+1:t_m+h}
-G_{d,h}H(S)x_{t_m}\right\|_2^2,
\end{equation}
where each historical validation block has ended
before the final origin. The smoothness decision is
\begin{equation}
\label{eq:intro_pooled_choice}
\widehat S^{\mathrm{FCV}}_{T,d,L,h}
\in\operatorname*{arg\,min}_{S\in[0,1]}
F^{\mathrm{pool}}_{T,d,L,h}(S).
\end{equation}
The final trend is refitted using the newest $L$
observations and extrapolated into an untouched
future test block. Thus the validation objective
is prospective: no future test observation is
available when the operational trend is chosen.

Normalized smoothness is an important computational
coordinate. It transforms the unbounded penalty
range $\lambda\in[0,\infty]$ into $S\in[0,1]$
without requiring an arbitrary upper cutoff.
At $S=0$ the historical observations are unsmoothed;
at $S=1$ the fit is projected onto the polynomial
null space of the difference penalty. Intermediate
values change the terminal differences of the fitted
trend and consequently the polynomial forecast
coefficients. For fixed difference order $d$, this
does not select the polynomial degree, which
remains at most $d-1$; it selects the smoothness
of the trend that is to be continued.

The ingredients themselves are established.
\citet{guerrero2007penalized,guerrero2008smoothness}
developed PLS and the percentage-of-smoothness
interpretation, while \citet{cortes2017smoothness}
examined attained smoothness under conventional
criteria in this journal. Controlled-smoothness
trends were already used for remittance and
exchange-rate forecasting
\citep{islas2019remittances,islas2025exchange}.
More directly, \citet{hart1994tscv} used
one-step-ahead predictive error to tune kernel
smoothing, and \citet{vilar2007nonparametric}
studied nonparametric time-series forecasting.
We therefore do not claim that predictive
smoothness tuning in general, or the index
itself, originated here. The statistical
object examined is their concrete combination
into horizon-matched \emph{future-block}
selection of normalized PLS trend smoothness.

A Monte Carlo study addresses whether matching
the validation horizon to the operational
forecast horizon changes smoothness and
out-of-sample error; how that choice compares
with ordinary CV, GCV, AICc, and
simulation-only recovery oracles; and how
the bounded optimization objective behaves
when it contains several minima. The frozen
study crosses five trend mechanisms, three
noise families, two noise levels, and 100
seeds, with six outer origins and four
forecast horizons, yielding 72,000 decisions.
These evaluations are repeated within seeds
and are not 72,000 independent replications.

The main result is an aggregate improvement
of horizon-matched tuning over one-step
tuning at longer horizons for this design.
Neither the mathematical reparameterization
nor the experiment establishes uniform
forecast dominance under all time-series
models. 
\section{Related work and contribution}
\label{sec:related_work}

\subsection{Penalized trends and established smoothing criteria}

Whittaker's graduation method \citep{whittaker1923}
and its PLS descendants provide the estimator
considered here. \citet{guerrero2007penalized}
gave a statistical formulation for trend PLS,
while \citet{guerrero2008smoothness} formalized
choosing a penalty from an analyst-specified
smoothness percentage. The normalized index
used below rescales this established trace-based
quantity to include the attainable full-range
limit one; it does not generate a new
smoothing operator.

The classical generalized cross-validation formulation of
\citet{craven1979gcv} underlies the GCV benchmark used
here. Of direct relevance to \emph{Communications in
Statistics---Simulation and Computation},
\citet{cortes2017smoothness} compared the
smoothness achieved by PLS under CV, GCV,
AICc, and BIC. Their simulation demonstrates
sensitivity of attained smoothness to
features of the underlying trend and sample
length. \citet{bates1987gcvpack}, also in
this journal, developed computational
routines for GCV. These studies establish
strong precedents for automatic regularization
tuning, but neither specifies the
$h$-step extrapolation MSE minimized
by \cref{eq:intro_pooled_loss}.

\subsection{Controlled-smoothness forecasting precedents}

Smoothing a series before forecasting is
not novel in itself. \citet{islas2019remittances}
used a controlled-smoothness trend inside
a multi-state Markov-switching model for
forecasting remittances to Mexico, and
\citet{islas2025exchange} used a related
framework for exchange-rate prediction.
The former reports a prescribed smoothness
percentage rather than direct selection of
that percentage from the forecast errors
of the simple finite-difference continuation.
These studies establish forecasting use of
controlled trends, while their downstream
predictive models and smoothing rules
are different from the criterion studied here.

The role of dependent noise in controlled-smoothness PLS
was examined by \citet{guerrero2018autocorr}; our
simulations include dependent disturbances but retain
an identity-weight smoothing operator.

\subsection{Forecast-based cross-validation}

The distinction between recovering a signal
and optimizing predictions is longstanding.
\citet{hart1994tscv} developed time-series
cross-validation for choosing the bandwidth
of a kernel smoother from one-step-ahead
predictions with dependent observations;
boundary extrapolation is relevant in
that construction. \citet{vilar2007nonparametric}
compared nonparametric time-series forecasting
methods and smoothing choices in this journal.
\citet{tashman2000oos} discusses rolling
forecast-origin evaluations, while
\citet{bergmeir2012cv,bergmeir2018validity}
study cross-validation and its validity
for time-series prediction in appropriate
settings. These are direct methodological
antecedents; our focus is an $h$-step
future-block loss for a finite-difference
PLS trend continued by a fixed operator.

\subsection{Recovery, likelihood, and extrapolation targets}

Recent PLS work illustrates why a smoothing
parameter should be defined together with
its loss target. \citet{franke2026hp}
calibrate HP smoothing using simulated
data with known growth trends and assess
\emph{latent trend recovery}. By contrast,
our deployable tuning rule uses only
historical realized forecast errors.
\citet{biessy2026whittaker} revisits
Whittaker--Henderson smoothing and studies
marginal-likelihood parameter selection
and constrained extrapolation, among
other questions. The endpoint problems
highlighted by \citet{hamilton2018hp}
remain relevant because a forecast
made by polynomial continuation may be
unstable even if smoothness was tuned
for future error.

\subsection{Direct comparison of objectives}

\begin{table}[t]
\centering
\small
\caption{Closest distinctions in how smoothness
is selected and evaluated.}
\label{tab:literature_positioning}
\begin{tabular}{@{}p{0.27\linewidth}p{0.26\linewidth}p{0.37\linewidth}@{}}
\toprule
Reference & Selection & Target\\
\midrule
\citet{guerrero2008smoothness} & User specifies $S$ & Trend regularity\\
\citet{cortes2017smoothness} & CV/GCV/AICc/BIC & Attained PLS smoothness\\
\citet{islas2019remittances} & Controlled smoothing & Markov-switching forecast\\
\citet{hart1994tscv} & Kernel bandwidth & One-step prediction\\
\citet{franke2026hp} & HP penalty & Latent trend recovery\\
This study & PLS $S\in[0,1]$ & Chronological $h$-step extrapolation MSE\\
\bottomrule
\end{tabular}
\end{table}

The particular methodological contribution
is a self-contained combination of the
normalized PLS smoothing coordinate,
declared finite-difference continuation,
horizon-matched chronological forecast
loss, and operational refitting.
The monotone $S$ transformation includes
both limiting trend models and supports
a compact numerical search, but it
does not alone alter the estimator,
remove multimodality, or prove better
predictive performance. The frozen
simulation evaluates the contribution
against alternatives with different
statistical targets.

\section{Finite-difference penalized trend estimation}

\subsection{Penalized least squares}

At forecast origin $T$, let
\begin{equation}
\label{eq:observation_window}
x_T=(y_{T-L+1},\ldots,y_T)^\top\in\mathbb{R}^{L}
\end{equation}
denote the estimation window. For difference order $d<L$, let
\begin{equation}
\label{eq:penalty_gram_matrix}
D_d\in\mathbb{R}^{(L-d)\times L},
\qquad
Q=D_d^\top D_d.
\end{equation}
The baseline estimator solves
\begin{equation}
\label{eq:pls_problem}
\widehat\tau_T(\lambda)
=
\arg\min_{\tau\in\mathbb{R}^{L}}
\left\{
\|x_T-\tau\|_2^2+\lambda\|D_d\tau\|_2^2
\right\},
\qquad \lambda\ge0,
\end{equation}
with solution
\begin{equation}
\label{eq:H}
\widehat\tau_T(\lambda)
=
H_\lambda x_T,
\qquad
H_\lambda=(I+\lambda Q)^{-1}.
\end{equation}

The paper uses the zero-drift, identity-weight formulation in \cref{eq:pls_problem} to isolate the smoothness-selection question. Guerrero's broader statistical formulation allows nonzero mean differences and more general covariance weighting \cite{guerrero2007penalized}; these are natural extensions but are not silently imposed in the baseline criterion.

\subsection{Spectral representation}

Because
\begin{equation}
\label{eq:penalty_positive_semidefinite}
v^\top Qv=\|D_dv\|_2^2\ge0,
\end{equation}
$Q$ is symmetric positive semidefinite. The dimension of its null
space, and the existence of the smoother at every finite penalty, follow
directly from the difference operator.

\begin{proposition}[Null space and finite-penalty well-posedness]
\label{prop:pls_wellposed}
For the standard consecutive order-$d$ difference operator with
$1\le d<L$,
\begin{equation}
\label{eq:penalty_rank}
\operatorname{rank}(D_d)=\operatorname{rank}(Q)=L-d,
\qquad
\ker(Q)=\ker(D_d),\qquad
\dim\ker(Q)=d.
\end{equation}
For every finite $\lambda\ge0$, $A_\lambda=I+\lambda Q$
and $H_\lambda=A_\lambda^{-1}$ are symmetric positive definite and
invertible. The objective in \cref{eq:pls_problem} therefore has the
unique minimizer $H_\lambda x_T$.
\end{proposition}

\begin{proof}
The condition $D_dv=0$ is the order-$d$ homogeneous
difference recurrence. Its first $d$ entries are arbitrary, and the
remaining entries are uniquely determined because the coefficient of
each successive new entry is nonzero. Thus
$\dim\ker(D_d)=d$, so rank--nullity gives
$\operatorname{rank}(D_d)=L-d$.
Clearly $D_dv=0$ implies $Qv=D_d^\top D_dv=0$.
Conversely, if $Qv=0$, then
$0=v^\top Qv=\|D_dv\|_2^2$, which implies $D_dv=0$.
The kernels are therefore equal, and rank--nullity also gives the
stated rank of $Q$. Finally, for every nonzero $v$,
\begin{equation}
\label{eq:finite_penalty_spd}
v^\top A_\lambda v
=\|v\|_2^2+\lambda\|D_dv\|_2^2>0
\qquad(0\le\lambda<\infty).
\end{equation}
Hence $A_\lambda$ is symmetric positive definite and invertible;
its inverse $H_\lambda$ is likewise positive definite. The Hessian
of the penalized objective is $2A_\lambda\succ0$; the objective is
strictly convex and coercive, so its stationary point $H_\lambda x_T$
is its unique minimizer.
\end{proof}

In particular, $Q$ has exactly $d$ zero eigenvalues and
$L-d$ strictly positive eigenvalues. Write
\begin{equation}
\label{eq:penalty_spectral_decomposition}
Q=
U\operatorname{diag}
\left(
\underbrace{0,\ldots,0}_{d},
\delta_1,\ldots,\delta_{L-d}
\right)U^\top,
\qquad
\delta_j>0.
\end{equation}
Then
\begin{equation}
\label{eq:H_spectral}
H_\lambda
=
U\operatorname{diag}
\left(
\underbrace{1,\ldots,1}_{d},
\frac{1}{1+\lambda\delta_1},
\ldots,
\frac{1}{1+\lambda\delta_{L-d}}
\right)U^\top.
\end{equation}

Thus the $d$ null-space directions are reproduced without shrinkage, while each penalizable direction is attenuated by
\begin{equation}
\label{eq:spectral_shrinkage}
\alpha_j(\lambda)=\frac{1}{1+\lambda\delta_j}.
\end{equation}
For $d=2$, the null space consists of discrete linear sequences, so increasing $\lambda$ moves the fitted path toward a linear trend.

For finite $\lambda$, the $d$ zero eigenvalues of $Q$ become
$d$ unit eigenvalues of $H_\lambda$; the remaining eigenvalues
belong to $(0,1]$. Consequently $\ker(H_\lambda)=\{0\}$ for
finite $\lambda$, despite $\dim\ker(Q)=d$. As $\lambda$ grows,
$A_\lambda$ can become ill-conditioned even though it remains
invertible, motivating spectral shrinkage and exact treatment
of the limiting projection rather than numerical inversion
at an arbitrarily large penalty.

\subsection{Normalized smoothness}

Guerrero's trace-based construction motivates measuring smoothing through the effective action of $H_\lambda$. We use
\begin{equation}
\label{eq:S}
S(\lambda)
=
1-
\frac{1}{L-d}
\sum_{j=1}^{L-d}
\frac{1}{1+\lambda\delta_j}.
\end{equation}

For comparison, the original Guerrero-type index is
$S^{\mathrm{G}}_d(\lambda)=1-\operatorname{tr}(H_\lambda)/L$.
Our full-range normalization is simply
\begin{equation}
\label{eq:smoothness_index_rescaling}
S(\lambda)=\frac{L}{L-d}S^{\mathrm{G}}_d(\lambda),
\qquad d\ge1.
\end{equation}
Thus the rescaling preserves the penalty ordering and every fitted trend;
it should not be interpreted as a new smoothing estimator
\cite{guerrero2007penalized,cortes2017smoothness}.

The normalization excludes the $d$ directions that cannot be penalized, so
\begin{equation}
\label{eq:smoothness_endpoints}
S(0)=0,
\qquad
\lim_{\lambda\to\infty}S(\lambda)=1.
\end{equation}

For a fixed linear smoother,
\begin{equation}
\label{eq:edf_definition}
\operatorname{edf}(\lambda)=\operatorname{tr}(H_\lambda).
\end{equation}
Combining \cref{eq:H_spectral} and \cref{eq:S},
\begin{equation}
\label{eq:edf}
\boxed{
\operatorname{edf}(\lambda)
=
L-(L-d)S(\lambda).
}
\end{equation}
Hence $S$ is a normalized complement of effective degrees of freedom.

The practical CV search uses $S$, not a truncated grid of $\lambda$.
The parameter transformation covers both limiting smoothing regimes:
$S=0$ corresponds to $\lambda=0$, whereas $S=1$ corresponds to
$\lambda\to\infty$, the least-squares projection onto the
$d$-dimensional null space of $D_d$. A uniform grid in $S$
therefore has an interpretable common bounded domain, unlike an
arbitrarily bounded grid in $\lambda$. Because the mapping is
strictly increasing, it does not invent new minima or make the
forecast loss convex; its advantage is coordinate choice and
endpoint handling.

\section{Forecast-optimal smoothness by chronological cross-validation}
\label{sec:forecast_cv}

\subsection{From a fitted trend to a prospective forecast}

The parameter of interest is the normalized smoothness $S$, not the
unbounded penalty $\lambda$. Define the continuous family of
smoothing matrices, including their limiting endpoint, by
\begin{equation}
\label{eq:sindexed_smoothing_matrix}
H(S)=
\begin{cases}
\bigl(I+\lambda(S)D_d^\top D_d\bigr)^{-1},&0\le S<1,\\
P_{\ker(D_d)},&S=1,
\end{cases}
\end{equation}
where $P_{\ker(D_d)}$ is the orthogonal projection onto the
null space of the finite-difference operator. Computing
$H(S)$ may require inverting the monotone map $S(\lambda)$,
but the domain on which cross-validation operates is the
\emph{full bounded interval} $[0,1]$.

A fitted trend is not automatically a forecast. We specify
a continuation rule that imposes zero future $d$th
differences, producing the fixed linear operator
\begin{equation}
\label{eq:continuation_operator_space}
G_{d,h}\in\mathbb{R}^{h\times L}.
\end{equation}
The corresponding prospective forecast from origin $t$ is
\begin{equation}
\label{eq:forecast_operator}
\widehat z_t(S)=G_{d,h}H(S)x_t.
\end{equation}
In terms of the final backward differences of the
historically fitted trend,
\begin{equation}
\label{eq:continuation}
\widehat\tau_{t+k\mid t}(S)
=\sum_{j=0}^{d-1}\binom{k+j-1}{j}
\nabla^j\bigl[H(S)x_t\bigr]_{L},
\qquad k=1,\ldots,h.
\end{equation}
For $d=1$, $d=2$, and $d=3$, these continuations are constant,
linear, and quadratic in the forecast step $k$, respectively.
More generally, fixing $d$ fixes the maximum degree $d-1$,
while changing $S$ changes the smoothed endpoint values and
differences, hence the coefficients of the forecast polynomial.
The method selects among \emph{future extrapolations}, not
merely among different post-hoc representations of the past.
Selecting $d$ itself is a separate model-selection decision;
it is held fixed within the baseline CV comparison.

\subsection{The forecast-MSE surface in the smoothness coordinate}

For an \emph{historical} forecast origin $t$, let
\begin{equation}
\label{eq:future_observation_block}
z_t=(y_{t+1},\ldots,y_{t+h})^\top
\end{equation}
be the subsequent realized validation block. The loss assigned
to a candidate smoothness $S\in[0,1]$ is
\begin{equation}
\label{eq:origin_loss_smoothness}
F_t(S)
=\frac1h\left\|z_t-G_{d,h}H(S)x_t\right\|_2^2.
\end{equation}
Expanding this expression makes the dependence on the
smoothness-indexed smoothing matrix explicit:
\begin{equation}
\label{eq:origin_loss_quadratic}
\begin{aligned}
h F_t(S)
={}&z_t^\top z_t
-2z_t^\top G_{d,h}H(S)x_t\\
&+x_t^\top H(S)G_{d,h}^\top
G_{d,h}H(S)x_t.
\end{aligned}
\end{equation}
The matrix $H(S)$ is symmetric. Unlike ordinary fitting
criteria, \cref{eq:origin_loss_smoothness} scores observations
\emph{after} the training endpoint; the in-sample residual
$\|x_t-H(S)x_t\|^2$ does not appear in the criterion.

For derivative calculations it is sometimes convenient to
use the algebraically equivalent auxiliary function
\begin{equation}
\label{eq:origin_loss}
f_t(\lambda)=F_t(S(\lambda)).
\end{equation}
This equality is a coordinate change, not a second
optimization problem or a requirement to search a truncated
range of $\lambda$.

\subsection{Chronological, horizon-matched forecast cross-validation}

At final forecast origin $T$, select historical origins
$t_1,\ldots,t_M$ whose training windows are available and
whose validation blocks are already observed:
\begin{equation}
\label{eq:chronological_origin_restriction}
t_m+h\le T,\qquad m=1,\ldots,M.
\end{equation}
For fixed $(d,L,h)$, the proposed CV surface is
\begin{equation}
\label{eq:F_S}
F^{\mathrm{pool}}_{T,d,L,h}(S)
=\frac1M\sum_{m=1}^M F_{t_m}(S),
\qquad S\in[0,1].
\end{equation}
It produces the forecast-optimal smoothness estimate
\begin{equation}
\label{eq:Sstar}
\widehat S^{\mathrm{FCV}}_{T,d,L,h}
\in\operatorname*{arg\,min}_{S\in[0,1]}
F^{\mathrm{pool}}_{T,d,L,h}(S).
\end{equation}
The selection space includes both exact limiting models,
$S=0$ and $S=1$, not just interior grid values.
When the objective has several minima, the global
minimizer is determined by the lowest historical predictive
MSE including the endpoints; a deterministic smallest-$S$
tie convention is used if several minima have the same
recorded objective value.

Both the forecasting horizon and the continuation operator
are essential parts of the criterion. Minimizing one-step
forecast error and deploying the resulting smoothness for
a longer horizon is a different procedure. The index makes
this comparison interpretable on $[0,1]$, while its
one-to-one transformation from $\lambda$ preserves the
fitted-trend family and any exact minimizing set.

\subsection{Final refit and untouched future forecast}

At outer origin $T$, all validation-stage trends are
discarded. Using the smoothness chosen from historical
blocks, the current trend is refitted on the newest
$L$ observations:
\begin{equation}
\label{eq:final_refit}
\widehat\tau_T
=H\bigl(\widehat S^{\mathrm{FCV}}_{T,d,L,h}\bigr)
y_{T-L+1:T}.
\end{equation}
The operational forecast is
\begin{equation}
\label{eq:final_operational_forecast}
\widehat y_{T+1:T+h\mid T}
=G_{d,h}\widehat\tau_T.
\end{equation}
No observation after $T$ enters either the CV decision
or the final fit. Historical forecast losses, not the
historical fitted trend paths, are averaged.

\subsection{Relationship to established cross-validation}

Ordinary leave-one-out CV and GCV primarily measure
interpolation or fitted-signal error for a smoother.
The present criterion uses chronological, genuinely
future-block MSE from a declared polynomial continuation.
This does not invent the general principle of predictive
cross-validation: Hart's time-series CV
\cite{hart1994tscv} is an important precedent.
The specific object investigated here is direct
$h$-step forecast-MSE selection for the
finite-difference PLS trend over normalized $S\in[0,1]$.
Time-adaptive decisions using historical minima, considered
later, are optional and not part of this definition.

\section{Basic properties}

\begin{proposition}[Monotone smoothness reparameterization]
\label{prop:monotone}
For positive penalized eigenvalues $\delta_1,\ldots,\delta_{L-d}$,
\begin{equation}
\label{eq:smoothness_first_derivative}
S'(\lambda)
=
\frac{1}{L-d}
\sum_{j=1}^{L-d}
\frac{\delta_j}{(1+\lambda\delta_j)^2}
>0,
\end{equation}
and
\begin{equation}
\label{eq:smoothness_second_derivative}
S''(\lambda)
=
-\frac{2}{L-d}
\sum_{j=1}^{L-d}
\frac{\delta_j^2}{(1+\lambda\delta_j)^3}
<0.
\end{equation}
After adjoining $\lambda=\infty$, $S$ maps $[0,\infty]$ continuously and strictly increasingly onto $[0,1]$.
\end{proposition}

\begin{proof}
Differentiate \cref{eq:S} term by term. Positivity of each $\delta_j$ gives the signs; the endpoint values follow from $(1+\lambda\delta_j)^{-1}$.
\end{proof}

The bounded scale has a common interpretation across fixed $(d,L)$ configurations while preserving the ordering of penalty strength.

\begin{proposition}[Effective degrees of freedom]
\label{prop:edf}
For fixed $\lambda$,
\begin{equation}
\label{eq:edf_smoothness_formula}
\operatorname{tr}(H_\lambda)=L-(L-d)S(\lambda).
\end{equation}
\end{proposition}

\begin{proof}
From \cref{eq:H_spectral},
\begin{equation}
\label{eq:edf_spectral_trace}
\operatorname{tr}(H_\lambda)
=
d+\sum_{j=1}^{L-d}(1+\lambda\delta_j)^{-1}.
\end{equation}
Substitution of \cref{eq:S} yields the result.
\end{proof}

\begin{proposition}[Equivalent optimization in penalty and smoothness space]
\label{prop:equiv}
Let $F(S)=f(\lambda(S))$. Then the global minimizers in $\lambda\in[0,\infty]$ and $S\in[0,1]$ correspond under the bijection in Proposition \cref{prop:monotone}. At interior points,
\begin{equation}
\label{eq:loss_change_of_variables}
F'(S)=\frac{f'(\lambda)}{S'(\lambda)},
\end{equation}
so interior stationary points correspond exactly.
\end{proposition}

\begin{proof}
A strictly increasing bijection preserves the minimizing set. The derivative identity follows from the chain rule.
\end{proof}

\begin{proposition}[Exact endpoints and the unique limiting fit]
\label{prop:endpoints}
Under the conditions of \cref{prop:pls_wellposed}, $H(0)=I$.
Let $U_0\in\mathbb R^{L\times d}$ have orthonormal columns spanning
$\ker(D_d)$. At $S=1$, interpreted as $\lambda\to\infty$,
\begin{equation}
\label{eq:exact_limiting_projection}
H(1)=H_\infty=U_0U_0^\top=P_{\ker(D_d)}.
\end{equation}
Although $H_\infty$ is singular with rank $d$ and nullity $L-d$,
the limiting trend is uniquely defined by
\begin{equation}
\label{eq:exact_limiting_constrained_fit}
H_\infty x_T
=\operatorname*{arg\,min}_{\tau:\,D_d\tau=0}
\|x_T-\tau\|_2^2.
\end{equation}
Thus no $1-\varepsilon$ surrogate is required for $S=1$.
\end{proposition}

\begin{proof}
At $\lambda=0$, $H_0=I$. By \cref{eq:H_spectral}, the
$d$ unit eigenvalues on $\ker(D_d)$ remain one, whereas
$(1+\lambda\delta_j)^{-1}\to0$ for every $\delta_j>0$.
This proves the projection formula. Any feasible
$\tau\in\ker(D_d)$ has the form $\tau=U_0a$.
Since $U_0^\top U_0=I$, the least-squares objective
$\|x_T-U_0a\|_2^2$ has the unique minimizer
$a=U_0^\top x_T$. Hence $\tau=U_0U_0^\top x_T$ is unique,
even though the projection matrix itself is not invertible.
\end{proof}

For $d=2$, the endpoints are the unsmoothed historical path and
its unique least-squares linear projection. The expression
$(I+\infty Q)^{-1}$ is not an ordinary matrix inverse:
$S=1$ must be computed as the exact spectral projection.
A large finite penalty or $S=1-\varepsilon$ is only an
approximation and can introduce avoidable numerical error.
Endpoint comparison is part of the scientific selection problem,
not merely a numerical detail.

\subsection{The future polynomial family}

The estimator produces a polynomial continuation only after the
future-difference rule has been specified. This matters for
interpreting forecast-CV as a prospective trend-selection method.

\begin{proposition}[Forecast polynomial indexed by smoothness]
\label{prop:continuation_polynomial}
Fix $1\le d<L$ and $h\ge1$. For every $S\in[0,1]$, the
continuation $k\mapsto\widehat\tau_{T+k\mid T}(S)$ is the
restriction to positive integer steps of a polynomial in
$k$ of degree at most $d-1$. The coefficients depend on
$S$ through the terminal backward differences of $H(S)x_T$.
\end{proposition}

\begin{proof}
For a fitted sequence $u=H(S)x_T$, imposing zero future
$d$th differences gives the discrete Newton continuation
\begin{equation}
\label{eq:newton_forecast_polynomial}
\widehat\tau_{T+k\mid T}(S)
=\sum_{j=0}^{d-1}\binom{k+j-1}{j}
(\nabla^j u)_L.
\end{equation}
Each binomial coefficient is a polynomial in $k$ of
degree $j$. The sum therefore has degree at most
$d-1$. Since $u=H(S)x_T$, its terminal differences
and consequently its coefficients vary with $S$.
The limiting case $S=1$ is well defined by the
orthogonal projection in \cref{prop:endpoints}.
\end{proof}

This proposition does not imply that the degree $d-1$
is chosen by forecast-CV: the baseline fixes $d$ and
selects among the corresponding extrapolations by
their historical future-block error. Changing $d$
would require an explicitly specified joint or
nested model-selection experiment.

\subsection{Existence and possible nonuniqueness}

The fitting problem has a unique solution for every $S\in[0,1]$: for $S<1$ this follows from strict convexity in \cref{prop:pls_wellposed}, and for $S=1$ from the constrained least-squares argument in \cref{prop:endpoints}. For finite $S<1$, $H_{\lambda(S)}$ depends continuously on $S$; the exact projection limit at $S=1$ makes the full family continuous on $[0,1]$. Thus, with a fixed finite collection of historical folds, $F^{\mathrm{pool}}_{T,d,L,h}(S)$ is continuous on the compact interval $[0,1]$ and attains at least one global minimum. \emph{Existence does not imply a unique minimizing $S$.}

Uniqueness is not assumed. Distinct smoothing regimes can produce similar validation loss, and the criterion may contain several local minima. This is statistically meaningful because competing minima can imply different effective degrees of freedom and terminal slopes. The implemented numerical search uses adaptive derivative bracketing, Brent root refinement, and exact endpoint comparison. The mathematical criterion in \cref{eq:Sstar} is independent of numerical search details and does not require tracking minima through time.

\subsection{Forecast optimality versus recovery optimality}

If a latent trend $\tau$ is known in simulation, a recovery-oriented oracle might choose
\begin{equation}
\label{eq:latent_recovery_oracle}
S_{\mathrm{rec}}
\in
\arg\min_S
\|\widehat\tau_S-\tau\|^2.
\end{equation}
The proposed criterion instead chooses $S^\star_{d,L,h}$ from future predictive error. These targets need not coincide. A smoother can reconstruct the historical latent path accurately but extrapolate poorly, or sacrifice historical recovery in exchange for more stable terminal derivatives.

This distinction is a central empirical question of the paper, not a claim that forecast-optimal smoothness is universally preferable under every loss.

\section{Simulation and evaluation design}
\label{sec:evaluation}

\subsection{Forecast-origin chronology and refitting}

The primary experiment asks whether a
chronological horizon-matched mean
squared forecast-error criterion gives
better smoothness decisions than
conventional smoothing criteria.
At each outer origin $T$, only completed
historical $h$-step validation blocks
can contribute to the selection:
$t_m+h\le T$. Every candidate $S$
is used to fit the window ending at
$t_m$, and the fitted trend is continued
through that historical validation
block. Once smoothness has been
selected, all intermediate fits are
discarded and a new trend is fitted
on the latest $L$ observations.
Neither the current forecast's future
outcomes nor its outer-test errors
are used for tuning.

The feasible comparators are
horizon-matched forecast-CV, one-step
forecast-CV, ordinary CV, generalized
CV (GCV), and AICc. Non-feasible
latent signal-recovery and latent
forecast oracles are computed only in
simulations and are distinguished
from operational selectors.

\subsection{Frozen factorial Monte Carlo design}

Checkpoint 03 fixes $d=2$ and $L=60$
and crosses five latent trend mechanisms,
three observation-noise mechanisms
(iid Gaussian, AR(1), and Student-$t$),
two noise scales, and 100 seeds.
The resulting $5\times3\times2\times100$
configurations comprise 3,000 scenarios.
Each scenario has six outer forecast origins
evaluated at $h\in\{1,3,6,12\}$, giving
\begin{equation}
\label{eq:simulation_sample_count}
3000\times6\times4=72{,}000
\end{equation}
outer-origin--horizon evaluations. Designs,
seeds, and output artifacts were frozen for
the paper-scale run. Innovations are reused
across some trend mechanisms within a seed,
so these evaluations are neither
statistically independent nor separate
confirmations of one preferred model.

\subsection{Losses and uncertainty units}

The tuning objective is historical
future-block mean squared error (MSE)
at the actual forecast horizon.
Outer-test accuracy is summarized by
root mean squared forecast error (RMSFE).
A geometric average of paired RMSFE
ratios below one favors the numerator
selection method. Such a geometric ratio
is not the same statistic as the
RMSFE obtained by pooling all test
errors over origins and horizons.
Comparisons use a common panel and
loss scale; their numbers should
not be pooled with distinct empirical
evaluation programs.

Seed-level resampling, rather than
treating every outer forecast as
independent, is appropriate for assessing
the variability of the paired
simulation comparisons. The checkpoint
provides descriptive seed-resampling
summaries. No uniform risk inequality
or universal statistical dominance
is inferred from their numerical values.

\subsection{Bounded smoothness optimization}

Numerical search takes place over
$S\in[0,1]$ and includes the
exact two limiting smoothing models.
Forecast-loss curves are not assumed
unimodal. Adaptive derivative sampling,
Brent refinement of bracketed roots,
and endpoint comparison produce
candidate minimizers without a
truncated penalty grid in $\lambda$.
These operations are a numerical
implementation of the statistical CV
criterion; the adaptive search is not
a mathematical certificate that every
possible stationary root was found.

\section{Monte Carlo evidence for forecast-optimal smoothness}
\label{sec:empirical}

The principal empirical question is whether choosing normalized PLS
smoothness by chronological, horizon-matched forecast MSE is useful
relative to established tuning criteria. The results below use the frozen pooled-selector
experiment, retained as historical evidence while final
numerical and simulation protocols are reassessed.

\subsection{Horizon-matched tuning and forecast error}

The pooled-selector simulation (Checkpoint 03) comprises 3,000 synthetic
scenarios, formed from five latent trends, three noise mechanisms, two noise
scales, and 100 random seeds. Each scenario has six chronological forecast
origins evaluated at horizons $h\in\{1,3,6,12\}$, giving 72,000
outer-origin--horizon evaluations. The estimator fixes $d=2$ and $L=60$.

\Cref{tab:cp03} reports the geometric root-mean-square forecast error
(RMSFE) ratio of horizon-matched pooled forecast-CV to two feasible selectors.
At $h=1$, horizon-matched and one-step forecast-CV coincide by construction.
At longer horizons the matched selector performs better in aggregate,
particularly under serially correlated disturbances. The results support the benefit of horizon matching within
this controlled design. They do not imply that every scenario
favors the proposed rule, or that the ranking holds under
all possible data-generating mechanisms. No branch tracking
is used in this experiment.

\begin{table}[t]
\centering
\caption{Frozen pooled-selector simulation (Checkpoint 03). Geometric RMSFE
ratios below one favor horizon-matched forecast-CV.}
\label{tab:cp03}
\begin{tabular}{@{}rcc@{}}
\toprule
$h$ & Relative to one-step FCV & Relative to GCV \\
\midrule
1  & 1.000 & 0.912 \\
3  & 0.947 & 0.816 \\
6  & 0.861 & 0.688 \\
12 & 0.762 & 0.554 \\
\bottomrule
\end{tabular}
\end{table}

Three figures visualize distinct features of the primary
forecast-CV experiment. \Cref{fig:cp03_loss_surfaces} illustrates
how the chosen horizon changes the forecast-MSE surface over
normalized smoothness. \Cref{fig:cp03_horizon_noise} compares
horizon-matched selection with one-step tuning across the noise
mechanisms. \Cref{fig:cp03_oracles} distinguishes the feasible
predictive choice from two simulation-only latent oracles.
The latent oracles are diagnostics, not deployable selection
rules.

\begin{figure}[t]
\centering
\includegraphics[width=\columnwidth]{figures/fig_sim05_objective_curves.pdf}
\caption{Checkpoint 03 example of horizon-specific forecast-loss
surfaces over normalized smoothness, with each curve scaled by its
minimum. The figure uses a fixed synthetic scenario with oscillatory
trend and AR(1) observation noise. It illustrates the effect of the
forecast horizon on the objective, not outer-test performance.}
\label{fig:cp03_loss_surfaces}
\end{figure}

\begin{figure}[t]
\centering
\includegraphics[width=\columnwidth]{figures/fig_sim01_horizon_matching_by_noise.pdf}
\caption{Checkpoint 03 geometric forecast-RMSFE ratios of
horizon-matched forecast-CV to one-step forecast-CV,
stratified by observation-noise mechanism. Intervals are
seed-resampling uncertainty summaries from the frozen
simulation; values below one favor horizon matching.}
\label{fig:cp03_horizon_noise}
\end{figure}

\begin{figure}[t]
\centering
\includegraphics[width=\columnwidth]{figures/fig_sim04_smoothness_targets.pdf}
\caption{Checkpoint 03 mean selected smoothness across forecast
horizons for feasible forecast-CV, a latent historical
trend-recovery oracle, and a latent future-forecast oracle.
Both oracles use unobserved simulation information and are
included solely to distinguish statistical targets.}
\label{fig:cp03_oracles}
\end{figure}


\subsection{Distinct recovery and prediction objectives}

The simulated latent-trend oracles displayed in
\cref{fig:cp03_oracles} clarify that historical
signal recovery and extrapolation are
different optimization targets. Both
oracles use unobserved latent information
and are not implementable alternatives to
historical forecast-CV. A smoother can
recover much of the historical trend yet
give unstable terminal differences when
continued into the future.

With difference order $d=2$, all compared
PLS forecasts use linear continuation.
The experiment holds this forecast class
fixed while varying the smoothness-selection
criterion; it does not establish the best
difference order or stochastic residual model.

\section{Interpretation and limitations}
\label{sec:discussion}

\subsection{A prediction-targeted smoothing choice}

Normalized smoothness describes the effective action
of the historical PLS smoother. For fixed $(d,L)$,
\begin{equation}
\label{eq:scope_smoothness_edf}
S=\frac{L-\operatorname{tr}\{H(S)\}}{L-d}.
\end{equation}
A value near zero preserves much of the historical
variation; a value near one approaches the
least-squares polynomial null space. When
the trend is to be forecast, the operative
criterion is the accuracy of its continuation,
not just its apparent historical regularity.
The same series may therefore warrant
different smoothness values at different
forecast horizons.

The continuation remains a discrete polynomial
of degree at most $d-1$ for fixed difference
order $d$. Changing $S$ changes the fitted
terminal level and differences and therefore
the coefficients of that future polynomial.
Forecast-CV does not, by itself, select the
polynomial degree. Joint order selection or
a stochastic residual forecast requires
a separately specified validation design.

\subsection{What the compact index does and does not do}

Using $S\in[0,1]$ avoids a finite arbitrary
upper boundary on $\lambda$ and gives the two
limiting smoothers an explicit role in
validation. However, $S(\lambda)$ is one-to-one:
the reparameterization preserves all fitted
trends and interior stationary points.
It is not a new smoothing estimator, it
does not guarantee a unique minimizer,
and it does not itself produce predictive
improvement. Large changes in the
unbounded penalty may be compressed
near $S=1$, so numerical search must
treat boundary regions carefully.

\subsection{Simulation scope and statistical caution}

The main evidence is a factorial Monte Carlo
comparison of smoothness \emph{selection rules}
for one specified PLS forecast-continuation
family. The paired decisions share seeds
and, in some conditions, innovations.
Thus the reported 72,000 outer evaluations
are not independent observations.
The aggregate geometric forecast-RMSFE
ratios do not establish uniform or
distribution-free optimality.

The baseline uses an identity-weight
PLS fit with equally spaced observations,
a fixed training length, fixed difference
order, and deterministic polynomial
continuation. AR(1) and heavy-tailed
disturbances enter the simulation, but
the fitting method does not estimate
their complete covariance structure.
Forecast-targeted smoothing cannot
eliminate all end-of-sample and
model-misspecification issues
\citep{hamilton2018hp}.
In particular, the method can perform
poorly if the true future evolution
is incompatible with its extrapolation
class.


\section{Conclusion}

Selecting smoothness for a trend that
will be extrapolated is a forecasting
problem, not merely a post-hoc
descriptive smoothing problem.
We formulated chronological,
horizon-matched forecast
cross-validation over a normalized
PLS smoothness index $S\in[0,1]$.
The criterion compares historical
future-block errors from an explicitly
defined continuation, includes the
two limiting smoothing regimes, and
refits the selected trend before
issuing a genuinely out-of-sample
forecast.

The PLS smoother, trace-based
smoothness index, and predictive
cross-validation principle are
established methods. The contribution
studied here is their particular
operational combination for
finite-difference trend extrapolation
and a frozen simulation comparison
with conventional smoothing criteria.
At horizons $3$, $6$, and $12$,
the simulated geometric RMSFE
ratios to one-step tuning were
$0.947$, $0.861$, and $0.762$,
respectively, supporting the value
of selecting smoothness for the
forecast horizon within this design.

The objective can have multiple
local minima, and changing from
$\lambda$ to $S$ does not remove them.
Consequently, practical selection
must compare competing numerical
candidates and both endpoints.
The evidence does not establish
universal superiority of forecast-CV,
optimality across all polynomial
continuation orders, or guarantees outside the specified loss and
finite-difference continuation family.
The principal result is a transparent
prediction-targeted tuning criterion
whose modeling assumptions,
computational search, and
information set can be stated
and evaluated explicitly.


\section*{Data availability}
The numerical experiments and their fixed simulation seeds are documented in the accompanying repository. The financial evaluation uses a frozen historical market-data snapshot with a machine-readable manifest; data availability and licensing remain subject to the original providers. Results using revised macroeconomic series are exploratory rather than claims of real-time-vintage forecasting performance.

\section*{Code availability}
Reusable implementations of finite-difference penalized trend estimation, chronological forecast objectives, rolling-origin validation, and smoothness transformations are maintained in the \texttt{trend\_estimation} package at \url{https://github.com/heri-espino/Trend-Estimation}. The repository also contains the manuscript source, experiment protocols, and frozen result artifacts used for the paper.

\section*{Conflict of interest}
The author declares no conflict of interest.

\bibliographystyle{plainnat}
\bibliography{references}

\end{document}
